\chapter{Dual Support Vector Machine}
\lecture{18}{4 Nov.}{}

\begin{prev}
    Linear SVM (\autoref{sec:svm}) is more robust and solvable with quadratic
    programming: the hard-margin problem is a QP in \((b, \bm{w})\), and
    feeding it transformed inputs \(\bm{z}_n = \bm{\Phi}(\bm{x}_n)\) gives a
    large-margin hyperplane with a sophisticated boundary. The catch, left
    at the end of the chapter, is that this QP has \(\tilde{d} + 1\)
    variables.
\end{prev}

This chapter removes that dependence on \(\tilde{d}\), or at least moves it
out of sight. It rewrites the SVM as an `equivalent' problem whose variables
are one number per \emph{example} rather than one per \emph{feature}, using
the Lagrange multipliers that regularization already hinted at. The new
problem is again a QP; its solution also tells us exactly which examples
the fattest hyperplane depends on.

\paragraph{Roadmap.}
\begin{enumerate}
    \item \textbf{Embedding Numerous Features: Kernel Models}
        \begin{itemize}
            \item[] Lecture 1: Linear Support Vector Machine --- linear SVM:
                more robust and solvable with quadratic programming
            \item[] \textbf{Lecture 2: Dual Support Vector Machine} ---
                Motivation of Dual SVM; Lagrange Dual SVM; Solving Dual SVM;
                Messages behind Dual SVM
        \end{itemize}
    \item Combining Predictive Features: Aggregation Models
    \item Distilling Implicit Features: Extraction Models
\end{enumerate}

\section{Motivation of Dual SVM}\label{sec:dual-motivation}

\source{slides p.2--6}

\subsection{Non-Linear Support Vector Machine Revisited}

The previous chapter solved the hard-margin SVM as a QP and ended by feeding
it transformed inputs. Written out, that non-linear SVM and the call to the
QP solver that computes it are:

\begin{center}
    \begin{minipage}[c]{0.4\textwidth}
        \[
            \begin{aligned}
                \min_{b, \bm{w}} \quad & \frac{1}{2} \bm{w}^\top \bm{w}\\
                \text{s.t.} \quad & y_n (\bm{w}^\top
                    \underbrace{\bm{z}_n}_{\bm{\Phi}(\bm{x}_n)} + b) \ge 1,\\
                & \text{for } n = 1, 2, \dots , N
            \end{aligned}
        \]
    \end{minipage}%
    \hfill
    \begin{minipage}[c]{0.56\textwidth}
        \begin{algorithm}[H]
            \caption{Non-Linear Hard-Margin SVM}\label{alg:nonlinear-hard-svm}
            \(\mathrm{Q} = \begin{bmatrix} 0 & \bm{0}_{\tilde{d}}^\top \\
                \bm{0}_{\tilde{d}} & \mathrm{I}_{\tilde{d}} \end{bmatrix}\);
                \(\bm{p} = \bm{0}_{\tilde{d}+1}\);
                \(\bm{a}_n^\top = y_n \begin{bmatrix} 1 & \bm{z}_n^\top
                \end{bmatrix}\); \(c_n = 1\)\;
            \(\begin{bmatrix} b \\ \bm{w} \end{bmatrix} \gets
                \mathrm{QP}(\mathrm{Q}, \bm{p}, \mathrm{A}, \bm{c})\)\;
            \Return \(b \in \mathbb{R}\) \& \(\bm{w} \in
                \mathbb{R}^{\tilde{d}}\) with \(g_{\mathrm{SVM}}(\bm{x}) =
                \sign(\bm{w}^\top \bm{\Phi}(\bm{x}) + b)\)\;
        \end{algorithm}
    \end{minipage}
\end{center}

This is \autoref{alg:linear-hard-svm} with \(\bm{x}_n\) replaced by
\(\bm{z}_n\) everywhere. It was built to serve two demands, and it also
inherits a cost:
\begin{itemize}
    \item What we demanded is a set of hypotheses that are \textbf{not many},
        thanks to the large margin, but have a \textbf{sophisticated}
        boundary, thanks to the feature transform --- the new possibility
        that closed \autoref{sec:why-large-margin}.
    \item What we pay is a QP with \(\tilde{d} + 1\) variables and \(N\)
        constraints, which is challenging if \(\tilde{d}\) is large,
        \textbf{or infinite?!} :-)
\end{itemize}
The two demands collide here. The more numerous the features, the better
the boundary, and the bigger the QP; a transform with infinitely many
features, which will turn out to be useful, cannot even be written down as a
QP of this form.

\begin{moral}
    Goal: SVM \textbf{without dependence on \(\tilde{d}\)}.
\end{moral}

\subsection{Todo: SVM `without' \texorpdfstring{\(\tilde{d}\)}{d~}}

The plan is to trade the original QP for another one of the same kind whose
size is counted by the number of examples instead:

\begin{center}
    \begin{minipage}[t]{0.45\textwidth}
        \textbf{Original SVM}\\
        (convex) QP of
        \begin{itemize}
            \item \(\tilde{d} + 1\) variables
            \item \(N\) constraints
        \end{itemize}
    \end{minipage}%
    \hfill
    \begin{minipage}[t]{0.45\textwidth}
        \textbf{`Equivalent' SVM}\\
        (convex) QP of
        \begin{itemize}
            \item \(N\) variables
            \item \(N + 1\) constraints
        \end{itemize}
    \end{minipage}
\end{center}

\begin{note}
    The aim is an SVM that stays computable even when \(\tilde{d}\) is huge.
    The QP in \autoref{alg:nonlinear-hard-svm} has \(\tilde{d} + 1\)
    variables, so its size grows with \(\tilde{d}\); the equivalent one is to
    be rewritten so that its variables and constraints are counted by \(N\)
    alone. Its cost then no longer depends on how many features
    \(\bm{\Phi}\) produces, so a large --- even infinite --- \(\tilde{d}\) is
    no obstacle.
\end{note}

\begin{remark}[Warning: Heavy Math!!!!!!]
    The road to the equivalent SVM runs through optimization theory, so two
    ground rules:
    \begin{itemize}
        \item We introduce some necessary math \emph{without rigor} to help
            understand SVM deeper.
        \item We `claim' some results if the details are unnecessary --- like
            how we `claimed' Hoeffding (\autoref{thm:hoeffding}).
    \end{itemize}
\end{remark}

\begin{moral}
    `Equivalent' SVM: based on some \textbf{dual problem} of Original SVM.
\end{moral}

\subsection{Key Tool: Lagrange Multipliers}

The tool comes from regularization, where the same move was already made
once, in \autoref{rem:lambda-c}.

\begin{center}
    \begin{tabular}{@{}c@{\quad}c@{\quad}c@{}}
        Regularization by & & Regularization by\\
        constrained-minimizing \(\Ein\) & \(\iff\) & minimizing
            \(\Eaug\)\\[4pt]
        \(\displaystyle \min_{\bm{w}} \Ein(\bm{w}) \text{ s.t. }
            \bm{w}^\top \bm{w} \le C\)
        & &
        \(\displaystyle \min_{\bm{w}} \Eaug(\bm{w}) = \Ein(\bm{w}) +
            \frac{\lambda}{N} \bm{w}^\top \bm{w}\)
    \end{tabular}
\end{center}

\begin{itemize}
    \item The constrained problem with budget \(C\) is equivalent to the
        augmented problem with some \(\lambda \ge 0\); the link is the
        \textbf{optimality condition}
        \[
            \nabla \Ein(\bm{w}) + \frac{2\lambda}{N} \bm{w} = \bm{0}.
        \]
    \item Regularization viewed \(\lambda\) as a \textbf{given parameter}
        instead of \(C\), which made the augmented problem `easy' to solve.
    \item The dual SVM goes the other way: it views the \(\lambda\)'s as
        unknown given the constraints, and \textbf{solves them as variables
        instead}.
\end{itemize}

In regularization the multiplier \(\lambda\) is the price that turns one
constraint into a penalty, and we never asked what its right value is for a
given \(C\) --- we just picked \(\lambda\). The SVM has its constraints
fixed, so the right multipliers are unknown, and finding them becomes the
problem itself.

\begin{moral}
    How many \(\lambda\)'s as variables? \(N\) --- one per constraint.
\end{moral}

\subsection{Starting Point: Constrained to `Unconstrained'}

Following that idea, we attach one multiplier to each of the SVM's \(N\)
constraints and fold the constraints into the objective, so that what
remains looks unconstrained.

\begin{center}
    \begin{minipage}[c]{0.38\textwidth}
        \[
            \begin{aligned}
                \min_{b, \bm{w}} \quad & \frac{1}{2} \bm{w}^\top \bm{w}\\
                \text{s.t.} \quad & y_n (\bm{w}^\top \bm{z}_n + b) \ge 1,\\
                & \text{for } n = 1, 2, \dots , N
            \end{aligned}
        \]
    \end{minipage}%
    \hfill
    \begin{minipage}[c]{0.58\textwidth}
        \begin{definition}[Lagrange function]
            With \vocab{Lagrange multipliers} \(\cancel{\lambda_n}\)
            \(\alpha_n\),
            \[
                \mathcal{L}(b, \bm{w}, \bm{\alpha}) =
                \underbrace{\frac{1}{2} \bm{w}^\top \bm{w}}_{\text{objective}}
                + \sum_{n=1}^N \alpha_n \underbrace{\left( 1 - y_n
                (\bm{w}^\top \bm{z}_n + b) \right)}_{\text{constraint}}.
            \]
        \end{definition}
    \end{minipage}
\end{center}

The SVM literature calls the multipliers \(\alpha_n\) rather than
\(\lambda_n\). Each constraint is rewritten as ``something \(\le 0\)'',
namely \(1 - y_n(\bm{w}^\top \bm{z}_n + b) \le 0\), and added to the objective
with its own multiplier.

\begin{lemma}[Claim]\label{lem:svm-minmax}
    \[
        \text{SVM} \equiv \min_{b, \bm{w}} \left( \max_{\text{all } \alpha_n
        \ge 0} \mathcal{L}(b, \bm{w}, \bm{\alpha}) \right)
        = \min_{b, \bm{w}} \left( \infty \text{ if violate};\
        \frac{1}{2} \bm{w}^\top \bm{w} \text{ if feasible} \right).
    \]
\end{lemma}

\begin{proof}
    Write \(\square = \frac{1}{2} \bm{w}^\top \bm{w}\) and fix \((b,
    \bm{w})\). The two words refer to the SVM's \(N\) constraints
    \(y_n (\bm{w}^\top \bm{z}_n + b) \ge 1\):
    \begin{itemize}
        \item \((b, \bm{w})\) is \vocab{feasible} if it satisfies \emph{all}
            of them, i.e.\ \(1 - y_n (\bm{w}^\top \bm{z}_n + b) \le 0\) for
            every \(n\) --- every example lies on its correct side, on or
            outside the margin;
        \item \((b, \bm{w})\) is \vocab{violating} if it fails \emph{at least
            one}, i.e.\ \(1 - y_n (\bm{w}^\top \bm{z}_n + b) > 0\) for some
            \(n\) --- that example is misclassified or sits inside the margin.
    \end{itemize}
    Every \((b, \bm{w})\) is exactly one of the two, so it suffices to compute
    the inner \(\max\) in each case.
    \begin{itemize}
        \item Any `violating' \((b, \bm{w})\): some constraint term \(1 - y_n
            (\bm{w}^\top \bm{z}_n + b)\) is positive, and
            \[
                \max_{\text{all } \alpha_n \ge 0} \left( \square + \sum_n
                \alpha_n (\text{some positive}) \right) \to \infty
            \]
            by sending that \(\alpha_n\) to infinity.
        \item Any `feasible' \((b, \bm{w})\): all constraint terms are
            non-positive, and
            \[
                \max_{\text{all } \alpha_n \ge 0} \left( \square + \sum_n
                \alpha_n (\text{all non-positive}) \right) = \square,
            \]
            attained by setting every \(\alpha_n = 0\).
    \end{itemize}
    The outer \(\min\) never picks an \(\infty\), so it ranges over the
    feasible \((b, \bm{w})\) only, and there minimizes \(\frac{1}{2}
    \bm{w}^\top \bm{w}\) --- which is the SVM.
\end{proof}

\begin{moral}
    Constraints now \textbf{hidden in} \(\max\).
\end{moral}

\begin{exercise}[Fun Time]
    Consider two transformed examples \((\bm{z}_1, +1)\) and \((\bm{z}_2,
    -1)\) with \(\bm{z}_1 = \bm{z}\) and \(\bm{z}_2 = -\bm{z}\). What is the
    Lagrange function \(\mathcal{L}(b, \bm{w}, \bm{\alpha})\) of hard-margin
    SVM?
    \begin{enumerate}[(1)]
        \item \(\frac{1}{2} \bm{w}^\top \bm{w} + \alpha_1 (1 + \bm{w}^\top
            \bm{z} + b) + \alpha_2 (1 + \bm{w}^\top \bm{z} + b)\);
        \item \(\frac{1}{2} \bm{w}^\top \bm{w} + \alpha_1 (1 - \bm{w}^\top
            \bm{z} - b) + \alpha_2 (1 - \bm{w}^\top \bm{z} + b)\);
        \item \(\frac{1}{2} \bm{w}^\top \bm{w} + \alpha_1 (1 + \bm{w}^\top
            \bm{z} + b) + \alpha_2 (1 + \bm{w}^\top \bm{z} - b)\);
        \item \(\frac{1}{2} \bm{w}^\top \bm{w} + \alpha_1 (1 - \bm{w}^\top
            \bm{z} - b) + \alpha_2 (1 - \bm{w}^\top \bm{z} - b)\).
    \end{enumerate}
\end{exercise}

\begin{proof}[Reference Answer]
    (2). By definition,
    \[
        \mathcal{L}(b, \bm{w}, \bm{\alpha}) = \frac{1}{2} \bm{w}^\top \bm{w}
        + \alpha_1 (1 - y_1 (\bm{w}^\top \bm{z}_1 + b)) + \alpha_2 (1 - y_2
        (\bm{w}^\top \bm{z}_2 + b))
    \]
    with \((\bm{z}_1, y_1) = (\bm{z}, +1)\) and \((\bm{z}_2, y_2) =
    (-\bm{z}, -1)\); the second term is \(1 + (-\bm{w}^\top \bm{z} + b) = 1
    - \bm{w}^\top \bm{z} + b\).
\end{proof}

\section{Lagrange Dual SVM}\label{sec:lagrange-dual}

\source{slides p.7--12}

\autoref{lem:svm-minmax} is still the original problem, only rewritten:
the \(\min\) is over \(\tilde{d} + 1\) variables. The step that removes them
is to swap the order of \(\min\) and \(\max\).

\subsection{Lagrange Dual Problem}

For any fixed \(\bm{\alpha}'\) with all \(\alpha'_n \ge 0\),
\[
    \min_{b, \bm{w}} \left( \max_{\text{all } \alpha_n \ge 0}
    \mathcal{L}(b, \bm{w}, \bm{\alpha}) \right) \ge \min_{b, \bm{w}}
    \mathcal{L}(b, \bm{w}, \bm{\alpha}'),
\]
because \(\max \ge\) any: for each \((b, \bm{w})\) the left integrand is at
least \(\mathcal{L}(b, \bm{w}, \bm{\alpha}')\), and taking \(\min\) on both
sides keeps the inequality.

For the best \(\bm{\alpha}' \ge \bm{0}\) on the right-hand side,
\[
    \min_{b, \bm{w}} \left( \max_{\text{all } \alpha_n \ge 0}
    \mathcal{L}(b, \bm{w}, \bm{\alpha}) \right) \ge
    \underbrace{\max_{\text{all } \alpha'_n \ge 0} \left( \min_{b, \bm{w}}
    \mathcal{L}(b, \bm{w}, \bm{\alpha}') \right)}_{\text{Lagrange dual
    problem}},
\]
because the best is one of any.

\begin{definition}[Lagrange dual problem]
    The \vocab{Lagrange dual problem} is the `outer' maximization over
    \(\bm{\alpha}\) of a lower bound on the original problem: each
    \(\bm{\alpha} \ge \bm{0}\) gives the lower bound \(\min_{b, \bm{w}}
    \mathcal{L}(b, \bm{w}, \bm{\alpha})\), and the dual looks for the
    tightest one.
\end{definition}

The swap is what makes the problem tractable: once \(\bm{\alpha}\) is
fixed, the inner \(\min\) over \((b, \bm{w})\) has \emph{no constraints}
left, so it can be solved by setting derivatives to zero.

\subsection{Strong Duality of Quadratic Programming}

A lower bound is useful only if it is tight. Placing the two problems side
by side,
\[
    \underbrace{\min_{b, \bm{w}} \left( \max_{\text{all } \alpha_n \ge 0}
    \mathcal{L}(b, \bm{w}, \bm{\alpha}) \right)}_{\text{equiv.\ to original
    (primal) SVM}}
    \ \ge \
    \underbrace{\max_{\text{all } \alpha_n \ge 0} \left( \min_{b, \bm{w}}
    \mathcal{L}(b, \bm{w}, \bm{\alpha}) \right)}_{\text{Lagrange dual}}
\]

\begin{definition}[Primal and dual]
    The \vocab{primal problem} is the problem we actually want to solve: the
    original hard-margin SVM, \(\min_{b, \bm{w}} \frac{1}{2} \bm{w}^\top
    \bm{w}\) s.t.\ \(y_n (\bm{w}^\top \bm{z}_n + b) \ge 1\), with variables
    \((b, \bm{w})\). By \autoref{lem:svm-minmax} it equals the left-hand
    side. The \vocab{dual problem} is the right-hand side, a new problem built
    from the primal whose variables are the multipliers \(\bm{\alpha}\).
\end{definition}

\begin{fact}[Strong duality of QP]\label{fact:strong-duality}
    Suppose the primal is a QP, i.e.\ its objective is quadratic and its
    constraints are linear, as the SVM is in
    \autoref{alg:nonlinear-hard-svm}.
    \begin{itemize}
        \item `\(\ge\)': \vocab{weak duality}, always true (shown above): the
            dual's optimal value never exceeds the primal's.
        \item `\(=\)': \vocab{strong duality}, the two optimal values
            coincide, so nothing is lost by solving the dual instead. For
            such a primal it holds as soon as
            \begin{itemize}
                \item the primal is convex (not automatic: a QP is convex
                    only if \(\mathrm{Q}\) is positive semi-definite),
                \item the primal is feasible, i.e.\ at least one \((b,
                    \bm{w})\) is feasible in the sense of the proof of
                    \autoref{lem:svm-minmax} (not automatic either: it depends
                    on the data; true if \(\bm{\Phi}\)-separable),
                \item the constraints are linear (automatic for a QP; it is
                    listed because the result extends to general convex
                    problems, where a non-linear constraint needs a point
                    satisfying it strictly, and plain feasibility suffices
                    only for linear ones)
            \end{itemize}
            --- called the \vocab{constraint qualification}. All three are
            conditions on the primal.
    \end{itemize}
\end{fact}

This is one of the results the warning promised to `claim'. The SVM meets
all three conditions: \(\frac{1}{2} \bm{w}^\top \bm{w}\) is convex, the
constraints are linear in \((b, \bm{w})\), and a hard-margin SVM is only
used on data that some hyperplane separates in \(\mathcal{Z}\).

\begin{remark}
    Strong duality is a statement about optimal \emph{values}: the minimum on
    the left and the maximum on the right are the same number. By itself it
    does not say that either side is reached by an actual point, since an
    infimum or a supremum need not be attained. A \vocab{primal-dual optimal}
    solution is a triple \((b, \bm{w}, \bm{\alpha})\) in which \((b,
    \bm{w})\) attains the primal optimum and \(\bm{\alpha}\) attains the dual
    optimum. Under strong duality such a triple exists exactly when
    \emph{both} sides are attained. Once they are, take any primal optimum
    \((b^*, \bm{w}^*)\) and any dual optimum \(\bm{\alpha}^*\), and write
    \(p^*\), \(d^*\) for the primal and dual optimal values. Then
    \[
        p^*
        \overset{\text{(i)}}{=}
        \max_{\bm{\alpha} \ge \bm{0}} \mathcal{L}(b^*, \bm{w}^*, \bm{\alpha})
        \overset{\text{(ii)}}{\ge}
        \mathcal{L}(b^*, \bm{w}^*, \bm{\alpha}^*)
        \overset{\text{(iii)}}{\ge}
        \min_{b, \bm{w}} \mathcal{L}(b, \bm{w}, \bm{\alpha}^*)
        \overset{\text{(iv)}}{=}
        d^*.
    \]
    (i) and (iv) say only that \((b^*, \bm{w}^*)\) attains the outer
    \(\min\) of the primal and \(\bm{\alpha}^*\) attains the outer \(\max\)
    of the dual. (ii) holds because a maximum over all \(\bm{\alpha} \ge
    \bm{0}\) is at least its value at the particular \(\bm{\alpha}^*\); (iii)
    holds because a minimum over all \((b, \bm{w})\) is at most its value at
    the particular \((b^*, \bm{w}^*)\). Strong duality makes the two ends
    equal, \(p^* = d^*\), so both inequalities in between are equalities.
    Thus the single triple \((b^*, \bm{w}^*, \bm{\alpha}^*)\) realizes both
    sides at once: \(\bm{\alpha}^*\) solves the primal's inner \(\max\) at
    \((b^*, \bm{w}^*)\), \((b^*, \bm{w}^*)\) solves the dual's inner
    \(\min\) at \(\bm{\alpha}^*\), and \(\mathcal{L}(b^*, \bm{w}^*,
    \bm{\alpha}^*)\) is the common optimal value. The two sides never ask
    for different points.

    For a QP both sides are attained automatically. A quadratic objective
    that is bounded on a non-empty set cut out by linear constraints always
    reaches its optimum (the Frank--Wolfe theorem). The primal SVM is bounded
    below by \(0\). The dual is also a quadratic program with linear
    constraints (\autoref{sec:solving-dual}), and by weak duality it is
    bounded above by the primal value.
\end{remark}

\begin{moral}
    There \textbf{exists a primal-dual optimal} solution \((b, \bm{w},
    \bm{\alpha})\) for both sides.
\end{moral}

So we may solve the dual instead of the primal, and the same triple
\((b, \bm{w}, \bm{\alpha})\) answers both.

\subsection{Solving Lagrange Dual: Simplifications (1/2)}

Strong duality lets us solve the dual instead of the primal. Written out in
full, with the primal variables \((b, \bm{w})\) and the dual variables
\(\bm{\alpha}\) told apart by colour, it reads:
\[
    \max_{\text{all } \dvar{\alpha_n} \ge 0} \left( \min_{\pvar{b},
    \pvar{\bm{w}}} \underbrace{\frac{1}{2} \pvar{\bm{w}}^\top \pvar{\bm{w}}
    + \sum_{n=1}^N \dvar{\alpha_n} \left( 1 - y_n (\pvar{\bm{w}}^\top
    \bm{z}_n + \pvar{b}) \right)}_{\mathcal{L}(\pvar{b}, \pvar{\bm{w}},
    \dvar{\bm{\alpha}})} \right)
\]

\begin{itemize}
    \item The inner problem is `unconstrained'; at the optimum
        \[
            \frac{\partial \mathcal{L}(\pvar{b}, \pvar{\bm{w}},
            \dvar{\bm{\alpha}})}{\partial \pvar{b}} = 0 = -\sum_{n=1}^N
            \dvar{\alpha_n} \cvar{y_n}.
        \]
    \item There is no loss of optimality if we solve with the constraint
        \(\cvar{\sum_{n=1}^N} \dvar{\alpha_n} \cvar{y_n = 0}\).
\end{itemize}

\begin{remark}
    The derivative is taken only in the inner problem. There \(\bm{\alpha}\)
    is held fixed, so \(\min_{b, \bm{w}} \mathcal{L}(b, \bm{w},
    \bm{\alpha})\) is an ordinary minimization of a differentiable function
    of \((b, \bm{w})\) over all of \(\mathbb{R}^{1 + \tilde{d}}\); this is
    what `unconstrained' means. Such a domain has no boundary, so a
    minimizer, if one exists, is an interior point and hence a critical
    point: every partial derivative vanishes there (Fermat's theorem). The
    outer \(\max\) is different. It runs over \(\bm{\alpha} \ge \bm{0}\),
    whose maximizer may lie on the boundary \(\alpha_n = 0\), so
    \(\partial \mathcal{L} / \partial \alpha_n = 0\) is never imposed.
\end{remark}

The second bullet deserves its reason. \(\mathcal{L}\) is linear in \(b\)
with slope \(-\sum_n \alpha_n y_n\). If that slope is non-zero, the inner
\(\min\) over \(b\) is \(-\infty\), and such an \(\bm{\alpha}\) can never be
the outer \(\max\); so the outer \(\max\) may as well be restricted to the
\(\bm{\alpha}\) with zero slope.

The new constraint has a pleasant side effect: \(\cvar{b}\)
\textbf{\cvar{can be removed}}. Under it the only term containing \(b\) is
\(-\sum_n \alpha_n y_n \cdot b = 0\), so the problem becomes
\[
    \max_{\text{all } \dvar{\alpha_n} \ge 0, \cvar{\sum y_n} \dvar{\alpha_n}
    \mathrel{\cvar{=}} \cvar{0}} \left( \min_{\pvar{b}, \pvar{\bm{w}}} \frac{1}{2}
    \pvar{\bm{w}}^\top \pvar{\bm{w}} + \sum_{n=1}^N \dvar{\alpha_n} \left( 1
    - y_n (\pvar{\bm{w}}^\top \bm{z}_n) \right) - \cancel{\cvar{\sum_{n=1}^N}
    \dvar{\alpha_n} \cvar{y_n} \cdot \pvar{b}} \right)
\]

\subsection{Solving Lagrange Dual: Simplifications (2/2)}

With \(b\) gone, the inner problem is a minimization over \(\bm{w}\) alone,
and the same trick applies to it:
\[
    \max_{\text{all } \dvar{\alpha_n} \ge 0, \cvar{\sum y_n} \dvar{\alpha_n}
    \mathrel{\cvar{=}} \cvar{0}} \left( \min_{\gone{b}, \pvar{\bm{w}}} \frac{1}{2}
    \pvar{\bm{w}}^\top \pvar{\bm{w}} + \sum_{n=1}^N \dvar{\alpha_n} \left( 1
    - y_n (\pvar{\bm{w}}^\top \bm{z}_n) \right) \right)
\]

\begin{itemize}
    \item The inner problem is `unconstrained'; at the optimum
        \[
            \frac{\partial \mathcal{L}(\pvar{b}, \pvar{\bm{w}},
            \dvar{\bm{\alpha}})}{\partial w_i} = 0 = \pvar{w_i} -
            \sum_{n=1}^N \dvar{\alpha_n} y_n z_{n,i}.
        \]
    \item There is no loss of optimality if we solve with the constraint
        \(\pvar{\bm{w}} \mathrel{\cvar{=}} \cvar{\sum_{n=1}^N} \dvar{\alpha_n} \cvar{y_n
        \bm{z}_n}\).
\end{itemize}

Here the inner objective is a convex quadratic in \(\bm{w}\), so the
stationary point is its minimum, and putting it in as a constraint just
records where the inner \(\min\) is attained.

This constraint also simplifies the objective. Under it,
\(\sum_n \alpha_n y_n \bm{w}^\top \bm{z}_n = \bm{w}^\top \sum_n \alpha_n y_n
\bm{z}_n = \bm{w}^\top \bm{w}\), so
\[
    \max_{\text{all } \dvar{\alpha_n} \ge 0, \cvar{\sum y_n}
    \dvar{\alpha_n} \mathrel{\cvar{=}} \cvar{0}, \pvar{\bm{w}} \mathrel{\cvar{=}} \cvar{\sum}
    \dvar{\alpha_n} \cvar{y_n \bm{z}_n}} \left( \gone{\min_{b, \bm{w}}}\
    \frac{1}{2} \pvar{\bm{w}}^\top \pvar{\bm{w}} + \sum_{n=1}^N
    \dvar{\alpha_n} - \pvar{\bm{w}}^\top \cvar{\bm{w}} \right)
\]
\[
    \iff \max_{\text{all } \dvar{\alpha_n} \ge 0, \cvar{\sum y_n}
    \dvar{\alpha_n} \mathrel{\cvar{=}} \cvar{0}, \pvar{\bm{w}} \mathrel{\cvar{=}} \cvar{\sum}
    \dvar{\alpha_n} \cvar{y_n \bm{z}_n}} -\frac{1}{2} \norm{\sum_{n=1}^N
    \dvar{\alpha_n} y_n \bm{z}_n}^2 + \sum_{n=1}^N \dvar{\alpha_n}.
\]
The \(\min\) is greyed out because nothing is left to minimize: \(b\) has
dropped out and \(\bm{w}\) is pinned down by \(\bm{\alpha}\). What remains is
a problem in \(\bm{\alpha}\) alone.

\subsection{KKT Optimality Conditions}

The dual is now a problem in \(\bm{\alpha}\) alone:
\[
    \max_{\text{all } \alpha_n \ge 0, \sum y_n \alpha_n = 0, \bm{w} = \sum
    \alpha_n y_n \bm{z}_n} -\frac{1}{2} \norm{\sum_{n=1}^N \alpha_n y_n
    \bm{z}_n}^2 + \sum_{n=1}^N \alpha_n
\]

Solving it gives the optimal \(\bm{\alpha}\), but what we finally need is
the hyperplane \((b, \bm{w})\). The bridge between the two is a list of
conditions that every primal-dual optimal triple must satisfy, collected from
the steps above.

\begin{fact}[KKT conditions]\label{fact:kkt}
    If \((b, \bm{w}, \bm{\alpha})\) is primal-dual optimal, then
    \begin{itemize}
        \item primal feasible: \(y_n (\bm{w}^\top \bm{z}_n + b) \ge 1\);
        \item dual feasible: \(\alpha_n \ge 0\);
        \item dual-inner optimal: \(\sum y_n \alpha_n = 0\); \(\bm{w} =
            \sum \alpha_n y_n \bm{z}_n\);
        \item primal-inner optimal (at optimal all `Lagrange terms'
            disappear):
            \[
                \alpha_n \left( 1 - y_n (\bm{w}^\top \bm{z}_n + b) \right) =
                0.
            \]
    \end{itemize}
    These are called the \vocab{Karush-Kuhn-Tucker (KKT) conditions},
    necessary for optimality [\& sufficient here].
\end{fact}

The first three restate constraints we have already met; the last is the
new one. It holds because \((b, \bm{w}, \bm{\alpha})\) solves the primal as
well as the dual, so \(\bm{\alpha}\) attains the primal's inner \(\max\)
over \(\bm{\alpha} \ge \bm{0}\) of \(\mathcal{L}(b, \bm{w}, \bm{\alpha})\)
(\autoref{lem:svm-minmax}). Each Lagrange term is
\(\alpha_n \ge 0\) times \(1 - y_n (\bm{w}^\top \bm{z}_n + b)\), and the
three possible signs of the latter go as in the proof of
\autoref{lem:svm-minmax}:
\begin{itemize}
    \item \(1 - y_n (\bm{w}^\top \bm{z}_n + b) > 0\) is impossible: sending
        \(\alpha_n \to \infty\) would make the inner \(\max\) infinite, and
        the outer \(\min\) never picks such \((b, \bm{w})\), because the
        primal is feasible (a condition of \autoref{fact:strong-duality}) and
        any feasible \((b, \bm{w})\) gives the finite value \(\frac{1}{2}
        \bm{w}^\top \bm{w}\);
    \item if \(1 - y_n (\bm{w}^\top \bm{z}_n + b) < 0\), the \(\max\) forces
        \(\alpha_n = 0\);
    \item if \(1 - y_n (\bm{w}^\top \bm{z}_n + b) = 0\), \(\alpha_n\) need
        not be \(0\), since the term is zero anyway.
\end{itemize}
In both remaining cases \(\alpha_n (1 - y_n (\bm{w}^\top \bm{z}_n + b)) = 0\).

As derived, the four conditions are only known to be necessary: every
primal-dual optimal \((b, \bm{w}, \bm{\alpha})\) satisfies them. For the SVM
they are also sufficient --- any \((b, \bm{w}, \bm{\alpha})\) satisfying all
four is primal-dual optimal --- because \(\mathcal{L}\) is convex in \((b,
\bm{w})\).

\begin{proof}[Proof of sufficiency]
    Let \((b, \bm{w}, \bm{\alpha})\) satisfy all four conditions. Write
    \(g(\bm{\alpha}') = \min_{b', \bm{w}'} \mathcal{L}(b', \bm{w}',
    \bm{\alpha}')\) for the dual objective, so the dual is \(\max_{\bm{\alpha}'
    \ge \bm{0}} g(\bm{\alpha}')\) and the primal is \(\min \frac{1}{2}
    \bm{w}'^\top \bm{w}'\) over feasible \((b', \bm{w}')\). The goal is to
    show that \((b, \bm{w})\) solves the primal and \(\bm{\alpha}\) solves
    the dual.

    \emph{Step 1: \(g(\bm{\alpha}) = \frac{1}{2} \bm{w}^\top \bm{w}\).}
    With \(\bm{\alpha}\) fixed, \(\mathcal{L}(\cdot, \cdot, \bm{\alpha})\) is
    a convex quadratic in \(\bm{w}\) plus a linear function of \(b\). The
    dual-inner conditions say exactly that its gradient vanishes at \((b,
    \bm{w})\), and for a convex function a zero gradient means a minimum, so
    \(g(\bm{\alpha}) = \mathcal{L}(b, \bm{w}, \bm{\alpha})\). Primal-inner
    optimality makes every Lagrange term zero, leaving
    \(\mathcal{L}(b, \bm{w}, \bm{\alpha}) = \frac{1}{2} \bm{w}^\top \bm{w}\).

    \emph{Step 2: \(g(\bm{\alpha}') \le \frac{1}{2} \bm{w}'^\top \bm{w}'\)
    for every \(\bm{\alpha}' \ge \bm{0}\) and every feasible \((b',
    \bm{w}')\).} This is weak duality, checked directly:
    \[
        g(\bm{\alpha}') \le \mathcal{L}(b', \bm{w}', \bm{\alpha}')
        = \frac{1}{2} \bm{w}'^\top \bm{w}' + \sum_{n=1}^N \alpha'_n \left(
        1 - y_n (\bm{w}'^\top \bm{z}_n + b') \right)
        \le \frac{1}{2} \bm{w}'^\top \bm{w}'.
    \]
    The first \(\le\) holds because a minimum over all \((b', \bm{w}')\) is
    at most the value at one of them; the second because each \(\alpha'_n
    \ge 0\) multiplies a factor that is \(\le 0\) at a feasible point.

    \emph{Step 3: \((b, \bm{w})\) solves the primal.} It is feasible by
    primal feasibility. Taking \(\bm{\alpha}' = \bm{\alpha}\) in Step~2 and
    using Step~1 gives \(\frac{1}{2} \bm{w}^\top \bm{w} \le \frac{1}{2}
    \bm{w}'^\top \bm{w}'\) for every feasible \((b', \bm{w}')\): no feasible
    point has a smaller objective.

    \emph{Step 4: \(\bm{\alpha}\) solves the dual.} It is feasible by dual
    feasibility. Taking \((b', \bm{w}') = (b, \bm{w})\) in Step~2, which is
    allowed because \((b, \bm{w})\) is feasible, and using Step~1 gives
    \(g(\bm{\alpha}') \le \frac{1}{2} \bm{w}^\top \bm{w} = g(\bm{\alpha})\)
    for every \(\bm{\alpha}' \ge \bm{0}\): no feasible \(\bm{\alpha}'\) has a
    larger dual objective.
\end{proof}

Convexity enters only in the first step. Without it a zero gradient need
not be a minimum, and the KKT conditions would be necessary but not
sufficient.

\begin{moral}
    We will use KKT to `solve' \((b, \bm{w})\) from the optimal
    \(\bm{\alpha}\).
\end{moral}

\begin{exercise}[Fun Time]
    For a single variable \(w\), consider minimizing \(\frac{1}{2} w^2\)
    subject to two linear constraints \(w \ge 1\) and \(w \le 3\). We know
    that the Lagrange function \(\mathcal{L}(w, \bm{\alpha}) = \frac{1}{2}
    w^2 + \alpha_1 (1 - w) + \alpha_2 (w - 3)\). Which of the following
    equations that contain \(\bm{\alpha}\) are among the KKT conditions of
    the optimization problem?
    \begin{enumerate}[(1)]
        \item \(\alpha_1 \ge 0\) and \(\alpha_2 \ge 0\);
        \item \(w = \alpha_1 - \alpha_2\);
        \item \(\alpha_1 (1 - w) = 0\) and \(\alpha_2 (w - 3) = 0\);
        \item all of the above.
    \end{enumerate}
\end{exercise}

\begin{proof}[Reference Answer]
    (4). (1) contains the dual-feasible constraints; (2) contains the
    dual-inner-optimal constraint, from \(\partial \mathcal{L} / \partial w
    = w - \alpha_1 + \alpha_2 = 0\); (3) contains the primal-inner-optimal
    constraints. (The solution is \(w = 1\), \(\alpha_1 = 1\), \(\alpha_2 =
    0\): the active constraint \(w \ge 1\) gets a positive multiplier, the
    inactive one gets zero.)
\end{proof}

\section{Solving Dual SVM}\label{sec:solving-dual}

\source{slides p.13--17}

\subsection{Dual Formulation of Support Vector Machine}

The previous section ended with the dual
\[
    \max_{\text{all } \alpha_n \ge 0, \sum y_n \alpha_n = 0, \bm{w} = \sum
    \alpha_n y_n \bm{z}_n} -\frac{1}{2} \norm{\sum_{n=1}^N \alpha_n y_n
    \bm{z}_n}^2 + \sum_{n=1}^N \alpha_n
\]

Flip the sign to turn \(\max\) into \(\min\), expand the squared norm as a
double sum, and drop \(\bm{w}\) from the constraints, since it no longer
appears in the objective and is recovered afterwards.

\begin{definition}[Standard hard-margin SVM dual]
    \[
        \begin{aligned}
            \min_{\bm{\alpha}} \quad & \frac{1}{2} \sum_{n=1}^N \sum_{m=1}^N
                \alpha_n \alpha_m y_n y_m \bm{z}_n^\top \bm{z}_m -
                \sum_{n=1}^N \alpha_n\\
            \text{subject to} \quad & \sum_{n=1}^N y_n \alpha_n = 0;\\
            & \alpha_n \ge 0, \text{ for } n = 1, 2, \dots , N.
        \end{aligned}
    \]
\end{definition}

\begin{moral}
    A (convex) QP of \(N\) variables \& \(N + 1\) constraints, as promised.
\end{moral}

Since it is a QP, the solver of \autoref{sec:svm} applies once its inputs
are written in the solver's format. :-)

\subsection{Dual SVM with QP Solver}

\begin{center}
    \begin{tabular}{@{}c@{\qquad}c@{}}
        optimal \(\bm{\alpha} = \ ?\) & optimal \(\bm{\alpha} \gets
            \mathrm{QP}(\mathrm{Q}, \bm{p}, \mathrm{A}, \bm{c})\)\\[6pt]
        \(\begin{aligned}
            \min_{\bm{\alpha}} \quad & \tfrac{1}{2} \textstyle\sum_{n=1}^N
                \sum_{m=1}^N \alpha_n \alpha_m y_n y_m \bm{z}_n^\top
                \bm{z}_m\\
            & \quad - \textstyle\sum_{n=1}^N \alpha_n\\
            \text{subject to} \quad & \textstyle\sum_{n=1}^N y_n \alpha_n =
                0;\\
            & \alpha_n \ge 0,\\
            & \text{for } n = 1, 2, \dots , N
        \end{aligned}\)
        &
        \(\begin{aligned}
            \min_{\bm{\alpha}} \quad & \tfrac{1}{2} \bm{\alpha}^\top
                \mathrm{Q} \bm{\alpha} + \bm{p}^\top \bm{\alpha}\\
            \text{subject to} \quad & \bm{a}_i^\top \bm{\alpha} \ge c_i,\\
            & \text{for } i = 1, 2, \dots
        \end{aligned}\)
    \end{tabular}
\end{center}

Matching the two forms term by term gives the solver's inputs:
\begin{itemize}
    \item the quadratic coefficients are \(q_{n,m} = y_n y_m \bm{z}_n^\top
        \bm{z}_m\);
    \item the linear coefficients are \(\bm{p} = -\bm{1}_N\);
    \item the constraint rows are \(\bm{a}_{\ge} = \bm{y}\) and
        \(\bm{a}_{\le} = -\bm{y}\) for the equality, and \(\bm{a}_n^\top =\)
        the \(n\)-th unit direction for each bound;
    \item the right-hand sides are \(c_{\ge} = 0\), \(c_{\le} = 0\) and
        \(c_n = 0\).
\end{itemize}

The generic solver only takes `\(\ge\)' constraints, so the one equality
\(\bm{y}^\top \bm{\alpha} = 0\) is written as the pair \(\bm{y}^\top
\bm{\alpha} \ge 0\) and \(-\bm{y}^\top \bm{\alpha} \ge 0\), and each bound
\(\alpha_n \ge 0\) as \(\bm{e}_n^\top \bm{\alpha} \ge 0\) --- \(N + 2\) rows
in all, which is the \(N + 1\) constraints with the equality counted twice.

\begin{note}
    Many solvers treat \textbf{equality} (\(\bm{a}_{\ge}\),
    \(\bm{a}_{\le}\)) \& \textbf{bound} (\(\bm{a}_n\)) constraints specially
    for \textbf{numerical stability}.
\end{note}

\subsection{Dual SVM with Special QP Solver}

A general-purpose solver works, but the dual has a structure that a
dedicated solver can exploit:
\[
    \text{optimal } \bm{\alpha} \gets \mathrm{QP}(\mathrm{Q}_{\mathrm{D}},
    \bm{p}, \mathrm{A}, \bm{c}): \qquad
    \begin{aligned}
        \min_{\bm{\alpha}} \quad & \frac{1}{2} \bm{\alpha}^\top
            \mathrm{Q}_{\mathrm{D}} \bm{\alpha} + \bm{p}^\top \bm{\alpha}\\
        \text{subject to} \quad & \text{special equality and bound
            constraints}
    \end{aligned}
\]

\begin{itemize}
    \item The coefficients \(q_{n,m} = y_n y_m \bm{z}_n^\top \bm{z}_m\) are
        often \textbf{non-zero}, so \(\mathrm{Q}_{\mathrm{D}}\) is dense.
    \item If \(N = 30{,}000\), the dense \(\mathrm{Q}_{\mathrm{D}}\) (\(N\)
        by \(N\) symmetric) takes \(> 3\)G RAM: that is \(9 \times 10^8\)
        entries, \(3.6\)\,GB even at four bytes each.
    \item To scale up to large \(N\), we need a special solver that
        \begin{itemize}
            \item does not store the whole \(\mathrm{Q}_{\mathrm{D}}\), and
            \item utilizes the special constraints properly.
        \end{itemize}
\end{itemize}

The subscript D marks the dual's matrix, to tell it from the \(\mathrm{Q}\)
of the primal. Its size is the new bottleneck: the dual traded
\(\tilde{d}\) for \(N\), and \(N^2\) entries is a lot when \(N\) is large.

\begin{moral}
    Usually better to use a \textbf{special solver} in practice.
\end{moral}

\subsection{Optimal \texorpdfstring{\((b, \bm{w})\)}{(b, w)}}

Whichever solver is used, it returns only \(\bm{\alpha}\). To turn it into
the hyperplane we return to the KKT conditions of \autoref{fact:kkt}, which
the optimal \(\bm{\alpha}\) satisfies together with the optimal \((b,
\bm{w})\): if \((b, \bm{w}, \bm{\alpha})\) is primal-dual optimal,
\begin{itemize}
    \item primal feasible: \(y_n (\bm{w}^\top \bm{z}_n + b) \ge 1\);
    \item dual feasible: \(\alpha_n \ge 0\);
    \item dual-inner optimal: \(\sum y_n \alpha_n = 0\); \(\bm{w} = \sum
        \alpha_n y_n \bm{z}_n\);
    \item primal-inner optimal (at optimal all `Lagrange terms' disappear):
        \[
            \alpha_n \left( 1 - y_n (\bm{w}^\top \bm{z}_n + b) \right) = 0
            \quad (\text{\vocab{complementary slackness}}).
        \]
\end{itemize}

Read as equations for \((b, \bm{w})\) given \(\bm{\alpha}\), they pin down
both:
\begin{itemize}
    \item The optimal \(\bm{w}\) is easy: the dual-inner condition gives it
        directly, \(\bm{w} = \sum \alpha_n y_n \bm{z}_n\).
    \item The optimal \(b\) takes one more step. Primal feasibility only
        gives a range for it, but complementary slackness gives an equality:
        if some \(\alpha_n > 0\), then \(b = y_n - \bm{w}^\top \bm{z}_n\).
\end{itemize}

The equality comes in one line. \(\alpha_n > 0\) forces \(1 - y_n
(\bm{w}^\top \bm{z}_n + b) = 0\); multiplying \(y_n (\bm{w}^\top \bm{z}_n +
b) = 1\) by \(y_n\), with \(y_n^2 = 1\), gives \(\bm{w}^\top \bm{z}_n + b =
y_n\). Some \(\alpha_n\) is always positive when both classes are present:
if all were zero, \(\bm{w} = \bm{0}\), and \(y_n b \ge 1\) cannot hold for
labels of both signs.

The same equation \(y_n (\bm{w}^\top \bm{z}_n + b) = 1\) has a geometric
reading. The optimal hyperplane separates the data, so the distance from
\(\bm{z}_n\) to it and the margin, the smallest such distance, are
\[
    \operatorname{distance}(\bm{z}_n, b, \bm{w}) = \frac{1}{\norm{\bm{w}}}
    y_n (\bm{w}^\top \bm{z}_n + b),
    \qquad
    \operatorname{margin}(b, \bm{w}) = \min_{m=1,\dots,N}
    \frac{1}{\norm{\bm{w}}} y_m (\bm{w}^\top \bm{z}_m + b) =
    \frac{1}{\norm{\bm{w}}},
\]
the last because at the optimum the closest example has \(y_m
(\bm{w}^\top \bm{z}_m + b) = 1\). An example with \(y_n (\bm{w}^\top
\bm{z}_n + b) = 1\) therefore sits at distance exactly
\(\frac{1}{\norm{\bm{w}}}\), the margin: it is one of the closest examples,
on the edge of the fat band, i.e.\ a support vector.

\begin{moral}
    Comp.\ slackness: \(\alpha_n > 0 \implies\) \(\bm{z}_n\) on fat
    boundary, i.e.\ \((\bm{z}_n, y_n)\) is a support vector (SV!).
\end{moral}

\begin{exercise}[Fun Time]
    Consider two transformed examples \((\bm{z}_1, +1)\) and \((\bm{z}_2,
    -1)\) with \(\bm{z}_1 = \bm{z}\) and \(\bm{z}_2 = -\bm{z}\). After solving
    the dual problem of hard-margin SVM, assume that the optimal
    \(\alpha_1\) and \(\alpha_2\) are both strictly positive. What is the
    optimal \(b\)?
    \begin{enumerate}[(1)]
        \item \(-1\);
        \item \(0\);
        \item \(1\);
        \item not certain with the descriptions above.
    \end{enumerate}
\end{exercise}

\begin{proof}[Reference Answer]
    (2). With the descriptions, at the optimal \((b, \bm{w})\),
    \[
        b = +1 - \bm{w}^\top \bm{z} = -1 + \bm{w}^\top \bm{z},
    \]
    the first from \((\bm{z}_1, +1)\) and the second from \((-\bm{z},
    -1)\). That is, \(\bm{w}^\top \bm{z} = 1\) and \(b = 0\).
\end{proof}

\section{Messages behind Dual SVM}\label{sec:dual-messages}

\source{slides p.18--22}

\subsection{Support Vectors Revisited}

Complementary slackness ties the dual solution to the picture of
\autoref{sec:svm}, where the examples on the edge of the fat band were first
called support vectors.

\begin{center}
    \begin{minipage}[c]{0.56\textwidth}
        \begin{itemize}
            \item The examples on the boundary `locate' the fattest
                hyperplane; the others are \textbf{not needed}.
            \item The examples with \(\alpha_n > 0\) lie on the boundary.
            \item We call the examples \((\bm{z}_n, y_n)\) with
                \(\alpha_n > 0\) the \textbf{support vectors}
                \cancel{(candidates)}; the boundary examples of
                \autoref{sec:svm} are only candidates.
            \item So SVs (positive \(\alpha_n\)) \(\subseteq\) SV candidates
                (on boundary).
        \end{itemize}
    \end{minipage}%
    \hfill
    \begin{minipage}[c]{0.4\textwidth}
        \centering
        \svmtoysolution
    \end{minipage}
\end{center}

The support vector \emph{candidates} are the examples on the boundary, \(y_n (\bm{w}^\top \bm{z}_n + b) = 1\). By
complementary slackness every example with \(\alpha_n > 0\) is a candidate.
The converse fails: for a candidate the Lagrange term \(\alpha_n (1 - y_n
(\bm{w}^\top \bm{z}_n + b))\) is zero whatever \(\alpha_n\) is, so the
solution may well give it \(\alpha_n = 0\). The candidates thus split into
those with \(\alpha_n > 0\) and those with \(\alpha_n = 0\), and the support
vectors are defined to be the former.

\begin{definition}[Support vector]
    An example \((\bm{z}_n, y_n)\) with \(\alpha_n > 0\) in the dual optimal
    solution is a \vocab{support vector} (SV).
\end{definition}

For the toy data of the picture (\(\bm{z}_n = \bm{x}_n\)),
the dual can be solved by hand from \(\bm{w} = (1, -1)\) and the KKT
conditions: \((3, 0)\) is off the boundary, so \(\alpha_4 = 0\), and
\[
    \bm{w} = -\alpha_2 (2, 2) + \alpha_3 (2, 0) = (1, -1),
    \qquad -\alpha_1 - \alpha_2 + \alpha_3 = 0
\]
give \(\alpha_2 = \frac{1}{2}\), \(\alpha_3 = 1\), \(\alpha_1 = \frac{1}{2}\).
Here all three candidates are SVs.

Since every non-SV has \(\alpha_n = 0\), the formulas for the optimal
hyperplane involve the SVs alone:
\begin{itemize}
    \item Only SVs are needed to compute \(\bm{w}\):
        \[
            \bm{w} = \sum_{n=1}^N \alpha_n y_n \bm{z}_n = \sum_{\text{SV}}
            \alpha_n y_n \bm{z}_n.
        \]
    \item Only SVs are needed to compute \(b\): \(b = y_n - \bm{w}^\top
        \bm{z}_n\) with any SV \((\bm{z}_n, y_n)\). In the toy example,
        \((2, 0)\) gives \(b = 1 - 2 = -1\).
\end{itemize}

\begin{moral}
    SVM: learn \textbf{fattest hyperplane} by identifying \textbf{support
    vectors} with the \textbf{dual} optimal solution.
\end{moral}

\subsection{Representation of Fattest Hyperplane}

The formula \(\bm{w} = \sum_n \alpha_n y_n \bm{z}_n\) has the same shape as
the weight vector of an algorithm from the very first chapters:

\begin{center}
    \begin{tabular}{@{}c@{\qquad\qquad}c@{}}
        \textbf{SVM} & \textbf{PLA}\\[4pt]
        \(\displaystyle \bm{w}_{\mathrm{SVM}} = \sum_{n=1}^N \alpha_n (y_n
            \bm{z}_n)\)
        &
        \(\displaystyle \bm{w}_{\mathrm{PLA}} = \sum_{n=1}^N \beta_n (y_n
            \bm{z}_n)\)\\[12pt]
        \(\alpha_n\) from \textbf{dual solution} & \(\beta_n\) by
            \textbf{\# mistake corrections}
    \end{tabular}
\end{center}

PLA (\autoref{alg:pla}) starts from \(\bm{0}\) and only ever adds \(y_n
\bm{z}_n\) for a mistaken example, so its final weight is that sum with
\(\beta_n\) counting how often example \(n\) was corrected.

\begin{moral}
    \(\bm{w}\) = linear combination of \(y_n \bm{z}_n\).
\end{moral}

\begin{itemize}
    \item The same holds for logistic or linear regression trained by GD/SGD
        from \(\bm{w}_0 = \bm{0}\): every gradient (\autoref{eq:lr-grad} for
        logistic regression) is itself a combination of the \(\bm{z}_n\), so
        every step adds one.
    \item Such a \(\bm{w}\) is said to be `\textbf{represented}' by the data.
\end{itemize}
What sets the SVM apart is how few terms the combination has.

\begin{moral}
    SVM: represent \(\bm{w}\) by \textbf{SVs only}.
\end{moral}

\subsection{Summary: Two Forms of Hard-Margin SVM}

We now have two QPs for the same hyperplane, each suited to a different
regime:

\begin{center}
    \begin{minipage}[t]{0.47\textwidth}
        \textbf{Primal Hard-Margin SVM}
        \[
            \begin{aligned}
                \min_{b, \bm{w}} \quad & \frac{1}{2} \bm{w}^\top \bm{w}\\
                \text{sub.\ to} \quad & y_n (\bm{w}^\top \bm{z}_n + b) \ge
                    1,\\
                & \text{for } n = 1, 2, \dots , N
            \end{aligned}
        \]
        \begin{itemize}
            \item \(\tilde{d} + 1\) variables, \(N\) constraints ---
                suitable when \(\tilde{d} + 1\) small
            \item physical meaning: locate \textbf{specially-scaled}
                \((b, \bm{w})\)
        \end{itemize}
    \end{minipage}%
    \hfill
    \begin{minipage}[t]{0.47\textwidth}
        \textbf{Dual Hard-Margin SVM}
        \[
            \begin{aligned}
                \min_{\bm{\alpha}} \quad & \frac{1}{2} \bm{\alpha}^\top
                    \mathrm{Q}_{\mathrm{D}} \bm{\alpha} - \bm{1}^\top
                    \bm{\alpha}\\
                \text{s.t.} \quad & \bm{y}^\top \bm{\alpha} = 0;\\
                & \alpha_n \ge 0 \text{ for } n = 1, \dots , N
            \end{aligned}
        \]
        \begin{itemize}
            \item \(N\) variables, \(N + 1\) simple constraints ---
                suitable when \(N\) small
            \item physical meaning: locate \textbf{SVs} \((\bm{z}_n, y_n)\)
                \& their \(\alpha_n\)
        \end{itemize}
    \end{minipage}
\end{center}

\begin{moral}
    Both eventually result in the optimal \((b, \bm{w})\) for the fattest
    hyperplane,
    \[
        g_{\mathrm{SVM}}(\bm{x}) = \sign(\bm{w}^\top \bm{\Phi}(\bm{x}) +
        b).
    \]
\end{moral}

\subsection{Are We Done Yet?}

The goal set at the start of the chapter was an SVM without dependence on
\(\tilde{d}\). The dual seems to deliver it:
\[
    \begin{aligned}
        \min_{\bm{\alpha}} \quad & \frac{1}{2} \bm{\alpha}^\top
            \mathrm{Q}_{\mathrm{D}} \bm{\alpha} - \bm{1}^\top \bm{\alpha}\\
        \text{subject to} \quad & \bm{y}^\top \bm{\alpha} = 0;\\
        & \alpha_n \ge 0, \text{ for } n = 1, 2, \dots , N
    \end{aligned}
\]
\begin{itemize}
    \item It has \(N\) variables and \(N + 1\) constraints, so its size shows
        no dependence on \(\tilde{d}\)?
    \item But each \(q_{n,m} = y_n y_m \bm{z}_n^\top \bm{z}_m\) is an inner
        product in \(\mathbb{R}^{\tilde{d}}\), which costs \(O(\tilde{d})\)
        via na\"ive computation!
\end{itemize}

The dependence has not gone away; it has moved into the entries of
\(\mathrm{Q}_{\mathrm{D}}\), each of which still transforms two inputs into
\(\mathbb{R}^{\tilde{d}}\) and multiplies them coordinate by coordinate.
The same holds at prediction time, where \(\bm{w}^\top \bm{\Phi}(\bm{x}) =
\sum_{\text{SV}} \alpha_n y_n \bm{z}_n^\top \bm{\Phi}(\bm{x})\) is again a
sum of inner products in \(\mathcal{Z}\).

\begin{moral}
    No dependence only if \textbf{avoiding na\"ive computation} (next
    lecture :-)).
\end{moral}

\begin{exercise}[Fun Time]
    Consider applying dual hard-margin SVM on \(N = 5566\) examples and
    getting \(1126\) SVs. Which of the following can be the number of
    examples that are on the fat boundary --- that is, SV candidates?
    \begin{enumerate}[(1)]
        \item \(0\);
        \item \(1024\);
        \item \(1234\);
        \item \(9999\).
    \end{enumerate}
\end{exercise}

\begin{proof}[Reference Answer]
    (3). Because SVs are always on the fat boundary,
    \[
        \#\text{SVs} \le \#\text{SV candidates} \le N,
    \]
    i.e.\ the answer must lie between \(1126\) and \(5566\).
\end{proof}

\hrulebar

\subsection*{Summary}

\begin{description}
    \item[Motivation of Dual SVM] want to remove dependence on
        \(\tilde{d}\).
    \item[Lagrange Dual SVM] KKT conditions link primal/dual.
    \item[Solving Dual SVM] another QP, better solved with a special
        solver.
    \item[Messages behind Dual SVM] SVs represent the fattest hyperplane.
\end{description}

Every place where \(\tilde{d}\) still enters the dual is an inner product
\(\bm{z}_n^\top \bm{z}_m\) of two transformed inputs. Next: computing the
inner product in \(\mathbb{R}^{\tilde{d}}\) efficiently, without forming the
transformed vectors at all.
